{
  \newlength{\LebarBersih}
  \setlength{\LebarBersih}{\dimexpr\textwidth-6\tabcolsep\relax}

  \newlength{\KolomDeskripsi}
  \newlength{\KolomIsi}

  \setlength{\KolomDeskripsi}{\dimexpr\LebarBersih/3\relax}
  \setlength{\KolomIsi}{\dimexpr\LebarBersih/3\relax}

  \newlength{\LebarTabel}
  \setlength{\LebarTabel}{%
    \dimexpr
      \KolomDeskripsi + 2\KolomIsi + 6\tabcolsep
    \relax
  }

  \newcommand{\JudulBlok}[4]{%
    \multicolumn{3}{@{}p{\LebarTabel}@{}}{%
      % \raggedright
      \vspace{-\dimexpr\baselineskip/2\relax}
      \hspace*{\tabcolsep}{%
        \hangindent=\tabcolsep
        \hangafter=1
        \rightskip=\tabcolsep
        #1\par
      }
      \vspace{-\dimexpr\baselineskip/2\relax}
      \rule{\LebarTabel}{0.4pt}\par
      % \vspace{-\dimexpr\baselineskip/8\relax}
      \begin{tabular}{%
        p{\KolomDeskripsi}%
        p{\KolomIsi}%
        p{\KolomIsi}}
        #2 & #3 & #4
      \end{tabular}
    }\\
  }

  \begin{longtable}{
    >{\raggedright\arraybackslash}p{\KolomDeskripsi}
    >{\raggedright\arraybackslash}p{\KolomIsi}
    >{\raggedright\arraybackslash}p{\KolomIsi}
  }
    \caption{Penelitian Terdahulu}\label{tbl:penelitianterdahulu}\\
    \toprule
    \textbf{Deskripsi} & \textbf{Kelebihan} & \textbf{Kekurangan} \\
    \midrule
    \endfirsthead

    \caption{Penelitian Terdahulu (lanjutan)} \\
    \toprule
    \textbf{Deskripsi} & \textbf{Kelebihan} & \textbf{Kekurangan} \\
    \midrule
    \endhead

    \midrule
    \multicolumn{3}{l}{\textit{Bersambung ke halaman berikutnya}}\\
    %\bottomrule
    \endfoot

    \bottomrule
    \endlastfoot

    % === Table Data ===
    \hline
    \JudulBlok{Audit Energi Sistem Kelistrikan Gedung Politeknik Negeri Ujung Pandang \cite{hamdani2017audit}}
    {Melakukan audit energi gedung dengan pengukuran parameter kelistrikan menggunakan \textit{power meter} dan sistem monitoring untuk menganalisis konsumsi energi dan kualitas data primer.}
    {Menghasilkan data pengukuran yang akurat dan mendukung identifikasi pemborosan energi secara detail.}
    {Proses audit dan analisis masih bersifat manual serta belum didukung oleh sistem perangkat lunak terintegrasi untuk pengolahan dan pelaporan.}
    \hline
    \hline
    \JudulBlok{Audit Energi Pada Gedung PPSDM Migas \cite{winarto2019audit}}
    {Menganalisis konsumsi energi gedung dengan fokus pada sistem tata udara dan pencahayaan melalui pencatatan durasi operasi, beban puncak, dan pola penggunaan energi.}
    {Memberikan gambaran praktis subsistem dominan konsumsi energi dan peluang konservasi yang signifikan.}
    {Belum menyediakan dukungan sistem informasi atau otomasi perhitungan dan dokumentasi audit.}
    \hline
    \hline
    \JudulBlok{Sistem Audit Intensitas Konsumsi Energi Listrik Berbasis GUI Matlab \cite{fitri2022sistem}}
    {Mengembangkan perangkat lunak berbasis antarmuka grafis untuk menghitung nilai IKE secara otomatis dari data konsumsi energi dan luas bangunan.}
    {Otomasi perhitungan IKE mengurangi kesalahan manual dan mempercepat proses evaluasi efisiensi energi.}
    {Cakupan sistem terbatas pada perhitungan IKE dan belum mendukung alur audit energi secara menyeluruh.}
    \hline
    \hline
    \JudulBlok{Evaluasi Intensitas Konsumsi Energi Melalui Audit Energi Listrik Gedung \cite{putra2025evaluasi}}
    {Menghitung IKE dan menggabungkannya dengan metode pengambilan keputusan multikriteria untuk menentukan prioritas strategi efisiensi energi.}
    {Menyediakan pendekatan terstruktur dalam penentuan prioritas intervensi berbasis data dan kriteria.}
    {Belum diimplementasikan dalam bentuk perangkat lunak audit energi yang terintegrasi.}
    \hline
    \hline
    \JudulBlok{Audit Support Systems and Decision Aids \cite{dowling2007audit}}
    {Mengkaji peran sistem pendukung audit dan tingkat restriktivitas sistem dalam mempengaruhi kualitas dan konsistensi proses audit.}
    {Memberikan kerangka konseptual desain sistem audit berbantuan perangkat lunak.}
    {Fokus pada audit laporan keuangan dan belum secara spesifik diterapkan pada audit energi bangunan.}
    \hline
  \end{longtable}
}